% !PRINTED 92
\weberpage{92}

%<<CONT
Determinante von $\varrho$ Reihen gebildet werden kann, die durch die
betreffende Primzahl nicht theilbar ist, und dass nicht alle Producte

\[
c_{1,h}\,e_{1,h}, \quad c_{2,h}\,e_{2,h}, \; \dots \; c_{\nu,h}\,e_{\nu,h}
\]

durch $q_h$ oder durch $q_h - 1$ theilbar werden. Dann setze man

\[
\gamma_{k,h} = c_{k,h}\,e_{k,h},
\]

und hat so nach (7) die Indices einer Basis einer Abel'schen Gruppe
$\mathfrak{B}$, deren reciproke Gruppe $\mathfrak{A}$ ein prim\"arer
Theiler der ganzen Gruppe $\mathfrak{N}$ ist.

\noindent\rule{4cm}{0.4pt}

\begin{center}
\S.\,22.

\gesperrt{Bestimmung der Gruppe $\mathfrak{A}$.}
\end{center}

Das letzte Ziel dieser Untersuchung ist nicht sowohl die Bestimmung der
Gruppe $\mathfrak{B}$, als die der Gruppe $\mathfrak{A}$, aus der man
direct die Kreistheilungsperioden $\eta = \Sigma r^a$ bilden kann. Diese
Aufgabe kann man auf folgende Art l\"osen: Nach (1) sind die Indices
$\alpha$ der Zahlen in $\mathfrak{A}$ von den $\beta$ abh\"angig durch
die Gleichung

\weq{1}{\omega_1^{\alpha_1\beta_1}\,\omega_2^{\alpha_2\beta_2} \dots
\omega_\mu^{\alpha_\mu\beta_\mu} = 1,}

worin $\omega_1, \omega_2, \dots \omega_\mu$ feste primitive
Einheitswurzeln der Grade $c_1, c_2, \dots c_\mu$ bedeuten. Versteht man
aber unter $C$ das kleinste gemeinschaftliche Vielfache von $c_1, c_2,
\dots c_\mu$ und setzt

\weq{2}{C = c_1 c_1' = c_2 c_2' = \dots = c_\mu c_\mu',}

und bezeichnet mit $\omega$ eine primitive $C^{\mathrm{te}}$
Einheitswurzel, so kann man statt (1) auch setzen

\[
\omega^{c_1'\alpha_1\beta_1 + c_2'\alpha_2\beta_2 + \dots +
c_\mu'\alpha_\mu\beta_\mu} = 1,
\]

und bekommt daher f\"ur die $\alpha$ die Congruenz

\[
c_1'\alpha_1\beta_1 + c_2'\alpha_2\beta_2 + \dots +
c_\mu'\alpha_\mu\beta_\mu \equiv 0 \pmod{C},
\]

die gleichbedeutend ist mit der Bedingung, dass die Summen

\[
\frac{\alpha_1\beta_1}{c_1} + \frac{\alpha_2\beta_2}{c_2} + \dots +
\frac{\alpha_\mu\beta_\mu}{c_\mu} = \sum_{1,\mu}^{h} \frac{\alpha_h
\beta_h}{c_h}
\]

ganze Zahlen sein m\"ussen. Wenn man darin f\"ur $\beta_k$ die
Ausdr\"ucke \S.\,21, (8) substituirt, so folgt, dass auch

\[
\sum^{k} x_k \sum^{h} \frac{\alpha_h\,\gamma_{k,h}}{c_h},
\]
%CONT>>
